% from ega2-1-fr.tex
\oldpage[II]{6}
noterons $\mathcal{A}(X)$ lorsqu'il n'en résultera pas de confusion~; si $U$
est un ouvert de $S$, on a
\[
  \mathcal{A}(f^{-1}(U))=\mathcal{A}(X)|U.
\]
De même, pour tout $\mathcal{O}_X$-Module $\mathcal{F}$ (resp. toute
$\mathcal{O}_X$-Algèbre $\mathcal{B}$), nous écrirons
$\mathcal{A}(\mathcal{F})$ (resp. $\mathcal{A}(\mathcal{B})$) l'image directe
$f_*(\mathcal{F})$ (resp. $f_*(\mathcal{B})$) qui est un
$\mathcal{A}(X)$-Module (resp. une $\mathcal{A}(X)$-Algèbre), et non seulement
un $\mathcal{O}_S$-Module (resp. $\mathcal{O}_S$-Algèbre).
\end{env}

\begin{env}[1.1.2]
\label{II.1.1.2-fr}
Soient $Y$ un second $S$-préschéma, $g:Y\to S$ son morphisme structural,
$h:X\to Y$ un $S$-morphisme~; on a donc le diagramme commutatif
\[
  \xymatrix{
    X\ar[rr]^h\ar[rd]_f && Y\ar[ld]^g\\
    & S
  }
  \tag{1.1.2.1}\label{II.1.1.2.1-fr}
\]

On a par définition $h=(\psi,\theta)$, où
$\theta:\mathcal{O}_Y\to h_*(\mathcal{O}_X)=\psi_*(\mathcal{O}_X)$ est un
homomorphisme de faisceaux d'anneaux~; on en déduit
(\hyperref[0.4.2.2-fr]{0, 4.2.2}) un homomorphisme de
$\mathcal{O}_S$-Algèbres
\[
  g_*(\theta):g_*(\mathcal{O}_Y)\to
  g_*(h_*(\mathcal{O}_X))=f_*(\mathcal{O}_X),
\]
autrement dit, un homomorphisme de $\mathcal{O}_S$-Algèbres
$\mathcal{A}(Y)\to\mathcal{A}(X)$, que nous noterons aussi
$\mathcal{A}(h)$. Si $h':Y\to Z$ est un second $S$-morphisme, il est immédiat
que $\mathcal{A}(h'\circ h)=\mathcal{A}(h)\circ\mathcal{A}(h')$. Nous avons
donc défini un \emph{foncteur contravariant} $\mathcal{A}(X)$ sur la catégorie
des $S$-préschémas, à valeurs dans la catégorie des
$\mathcal{O}_S$-Algèbres.

Soient maintenant $\mathcal{F}$ un $\mathcal{O}_X$-Module, $\mathcal{G}$ un
$\mathcal{O}_Y$-Module, et $u:\mathcal{G}\to\mathcal{F}$ un $h$-morphisme,
c'est-à-dire (\hyperref[0.4.4.1-fr]{0, 4.4.1}) un homomorphisme de
$\mathcal{O}_Y$-Modules $\mathcal{G}\to h_*(\mathcal{F})$. Alors
\[
  g_*(u):g_*(\mathcal{G})\to g_*(h_*(\mathcal{F}))=f_*(\mathcal{F})
\]
est un homomorphisme $\mathcal{A}(\mathcal{G})\to
\mathcal{A}(\mathcal{F})$ de $\mathcal{O}_S$-Modules, que nous noterons
$\mathcal{A}(u)$~; en outre, le couple
$(\mathcal{A}(h),\mathcal{A}(u))$ constitue un \emph{di-homomorphisme} du
$\mathcal{A}(Y)$-Module $\mathcal{A}(\mathcal{G})$ dans le
$\mathcal{A}(X)$-Module $\mathcal{A}(\mathcal{F})$.
\end{env}

\begin{env}[1.1.3]
\label{II.1.1.3-fr}
Le préschéma $S$ étant fixé, on peut considérer les couples
$(X,\mathcal{F})$, où $X$ est un $S$-préschéma et $\mathcal{F}$ un
$\mathcal{O}_X$-Module, comme formant une \emph{catégorie}, en définissant un
\emph{morphisme} $(X,\mathcal{F})\to(Y,\mathcal{G})$ comme un couple $(h,u)$,
où $h:X\to Y$ est un $S$-morphisme et $u:\mathcal{G}\to\mathcal{F}$ un
$h$-morphisme. On peut alors dire que
$(\mathcal{A}(X),\mathcal{A}(\mathcal{F}))$ est un \emph{foncteur
contravariant} à valeurs dans la catégorie dont les objets sont les couples
formés d'une $\mathcal{O}_S$-Algèbre et d'un Module sur cette Algèbre, et les
morphismes sont les di-homomorphismes.
\end{env}

\subsection{Préschémas affines sur un préschéma}
\label{subsection:II.1.2-fr}

\begin{definition}[1.2.1]
\label{II.1.2.1-fr}
Soient $X$ un $S$-préschéma, $f:X\to S$ son morphisme structural. On dit que
$X$ est \emph{affine au-dessus de $S$} s'il existe un recouvrement
$(S_\alpha)$ de $S$ par des ouverts affines tels que, pour tout $\alpha$, le
préschéma induit par $X$ sur l'ensemble ouvert $f^{-1}(S_\alpha)$ soit affine.
\end{definition}

\begin{example}[1.2.2]
\label{II.1.2.2-fr}
Tout sous-préschéma fermé de $S$ est un $S$-préschéma affine au-dessus de $S$
(\hyperref[I.4.2.3-fr]{I, 4.2.3} et
\hyperref[I.4.2.4-fr]{4.2.4}).
\end{example}

\begin{remark}[1.2.3]
\label{II.1.2.3-fr}
Un préschéma $X$ affine au-dessus de $S$ n'est pas nécessairement un schéma
affine, comme le montre l'exemple $X=S$ (\hyperref[II.1.2.2-fr]{1.2.2}).
D'autre part, si un schéma affine $X$ est un $S$-préschéma, $X$ n'est pas
nécessairement affine au-dessus
